\section{The linear program and the running instance}\label{sec:model}

Let $\Gset=(\Vset,\Aset)$ be a directed acyclic graph (DAG), with
$n=|\Vset|$ vertices and $m=|\Aset|$ arcs. An arc $(i,j)$ means that
\emph{vertex $j$ is a prerequisite of vertex $i$}. Its inequality is therefore
$x_i\le x_j$. This convention is used in every formula and diagram below.
Each vertex has profit $p_i\in\Rset$ and weight $w_i>0$. There is a single vertex
type: no bipartition or alternation of vertex types is assumed, a vertex may have
prerequisites and itself be a prerequisite, and profits may be negative, so
selecting a profitable vertex can require selecting loss-making predecessors.

\begin{definition}[Closure]
A set $C\subseteq\Vset$ is a closure if $i\in C$ and $(i,j)\in\Aset$
imply $j\in C$. Write $\Cset$ for all closures, including the empty set.
For $S\subseteq\Vset$, let
\[
 p(S)=\sum_{i\in S}p_i,\qquad w(S)=\sum_{i\in S}w_i,
 \qquad \varrho(S)=\frac{p(S)}{w(S)}\quad(S\ne\emptyset).
\]
\end{definition}

\subsection{Names and notation, taken from the literature}\label{sec:names}
Every term used below already has an owner, and this tutorial uses the owner's
word rather than a new one.

\begin{center}
\begin{tabular}{@{}lll@{}}
\toprule
Object & Name used here & Source \\
\midrule
$\Gset=(\Vset,\Aset)$ & digraph, vertices, arcs & standard \\
$C\subseteq\Vset$ down-closed & \emph{closure} & \citet{Picard1976} \\
$\max_C p(C)/w(C)$ & \emph{maximum-ratio closure} & \citet{Hochbaum2024} \\
the same set, order-theoretic & $\rho$-maximal \emph{initial set} & \citet{Sidney1975,Lawler1978} \\
the nested chain of closures & \emph{Sidney decomposition} & \citet{Sidney1975} \\
its successive differences $\Kset_r$ & \emph{blocks} & this tutorial \\
forest with one distinguished vertex per tree & \emph{rooted} forest, \emph{roots} & \citet{Yang2026} \\
the tree with no such vertex & \emph{rootless} tree & \citet{Yang2026} \\
forest with per-branch aggregates & \emph{normalized tree} & \citet{LerchsGrossmann1965} \\
\bottomrule
\end{tabular}
\end{center}

\noindent
\citet{Picard1976} writes \emph{maximal} closure where current usage writes
\emph{maximum} closure; the latter is used here. Matrices and vectors are set in
bold, scalars and indexed entries are not, so $\vect{w}$ is the weight vector
and $w_i$ one of its entries.

\subsection{The problem: maximum-ratio closure}\label{sec:problem}
The \FS{} solves one problem, and it is a \emph{ratio} problem over closures:
\begin{equation}\label{eq:ratioproblem}
  \rho^* \;=\; \max\Bigl\{\,\varrho(C)=\frac{p(C)}{w(C)}
              \;:\; \emptyset\ne C\subseteq\Vset,\ C\in\Cset \,\Bigr\}.
\end{equation}
Note what is absent: there is no capacity, and no upper bound on any variable.
The name of the method follows \eqref{eq:ratioproblem}, in the way that
\emph{network simplex} names the class of flow problems rather than the spanning
tree it maintains.

Two devices linearize \eqref{eq:ratioproblem}. Normalizing the denominator
\citep{CharnesCooper1962} gives a single linear program, stated as
\eqref{eq:normalized} and proved equivalent in \cref{thm:normalization}; that
program, and its structure, is what the rest of this tutorial works on, and
\cref{sec:class} identifies the class it belongs to. Treating the ratio
parametrically instead \citep{Dinkelbach1967} turns each iterate into a maximum
closure and hence a minimum cut \citep{Picard1976}; that is the route of the
comparison methods in \cref{app:flow}, not of the method developed here.

\subsection{Application: adding a capacity row}\label{sec:capacity}
The program \eqref{eq:ratioproblem} is an oracle, and the problem it serves is the
LP relaxation of a precedence-constrained knapsack. For capacity $c\ge0$,
\begin{subequations}\label{eq:lp}
\begin{align}
 z(c)=\max\quad &\sum_{i\in\Vset}p_i x_i,\label{eq:lpobj}\\
 \text{s.t.}\quad &\sum_{i\in\Vset}w_i x_i\le c,\label{eq:lpcap}\\
 &x_i-x_j\le0 &&(i,j)\in\Aset,\label{eq:lpprec}\\
 &0\le x_i\le1 &&i\in\Vset.\label{eq:lpbox}
\end{align}
\end{subequations}
The indicator $\ind^C$ is feasible for the precedence and box constraints
exactly when $C$ is a closure; the capacity row makes the LP generally
fractional. Solving \eqref{eq:lp} for a given $c$ needs no new machinery beyond
\eqref{eq:ratioproblem}: repeated calls of the oracle on residual graphs produce
closures of decreasing ratio, and $c$ decides how far along that sequence one
walks and which fraction of the critical block is taken.
\Cref{sec:outer} does this, and \cref{thm:outer} gives the whole function
$z(\cdot)$ at once, piecewise linear with one breakpoint per block, at no
extra cost. The capacity therefore never appears inside the simplex; it appears
only in the procedure that calls it, which is why the method is named after
\eqref{eq:ratioproblem} and not after the knapsack.

The same machinery answers a question in which no capacity appears at all.
Driving the oracle to exhaustion --- extract a block, delete it, ask again on
the residual until nothing is left --- produces the \emph{canonical
decomposition}: the complete sequence $\Kset_1,\dots,\Kset_k$ of blocks of
strictly decreasing ratio, together with the threshold of every vertex and the
breakpoints of the parametric closure value. This is the object the method is
built to deliver, and it is what the solver computes under the tasks
\texttt{full-path} and \texttt{positive-path}; a single capacity is then just a
query against a sequence already in hand. The sequence is not a device of this
tutorial: it is the decomposition of \citet{Sidney1975}, and the same object is
studied on precedence forests, under the name \emph{closure layers}, by the
companion work of this project. What is specific here is that a simplex basis,
rather than a flow, is what walks it.

\subsection{Eight vertices, several precedence levels}
We use the running instance from the supplied manuscript. Its data are
stored once in \texttt{tests/fixtures/paper\_example8.json}; the tables and
figures below are generated from that file.
\begin{table}[htbp]
\centering\input{../build/tutorial/generated/instance_table.tex}
\caption{Profits and weights of the running instance.}\label{tab:data}
\end{table}
Its arcs, in the fixed order used to number the slack variables, are
\begin{equation}\label{eq:arcs}
\begin{split}
\Aset=\{&(5,1),(5,2),(5,6),(6,3),(7,3),\\
         &(7,4),(8,5),(8,6),(8,7)\}.
\end{split}
\end{equation}
For example, vertex 8 requires vertices 5, 6 and 7, which in turn impose other
prerequisites. At $c=4$, the closure $\{1,3,6\}$ is an optimal integer
selection of value 5. The LP value will be 7. The integer problem is used
only to illustrate the difference; it is not the problem solved here.
\begin{figure}[htbp]
\centering\input{../build/tutorial/generated/graph.tex}
\caption{The general precedence graph. A label $(p_i,w_i)$ is attached to
each vertex; arrows point to prerequisites.}\label{fig:graph}
\end{figure}

\subsection{The dual and a convention to remember}
For arc multipliers $\alpha$, use the divergence convention of
\citet{DoseFuriniLocatelli2026}:
\[
\divg_i(\alpha)=\sum_{j:(i,j)\in\Aset}\alpha_{ij}
                  -\sum_{j:(j,i)\in\Aset}\alpha_{ji}.
\]
The dual of \eqref{eq:lp} is
\begin{equation}\label{eq:lpdual}
\begin{aligned}
\min\quad &c\lambda+\sum_i\mu_i\\
\text{s.t.}\quad &w_i\lambda+\mu_i+\divg_i(\alpha)\ge p_i &&i\in\Vset,\\
&\lambda\ge0,\quad\mu_i\ge0,\quad\alpha_{ij}\ge0.
\end{aligned}
\end{equation}
The multiplier $\lambda$ prices capacity; $\mu_i$ prices the upper bound.
Later, $s_{ij}$ will denote a \emph{primal slack}, not an arc multiplier.
Keeping $s_{ij}$, $\alpha_{ij}$ and $\lambda$ distinct prevents a collision
with the notation in the historical forest notes.

\subsection{Where this sits in the literature}\label{sec:prior}
Nothing in \eqref{eq:ratioproblem} or in the forest structure that follows from it is
new, and it is better to say so here than to let a reader discover it later.
The companion document in \path{literature/} does this carefully; the four
blocks are:

\begin{description}[leftmargin=*,style=nextline]
\item[Closure and its reduction to flow.] Maximum closure and its solution by
  minimum cut are classical \citep{Picard1976,PicardQueyranne1982}; the
  pseudoflow algorithm \citep{Hochbaum2008} is what a flow-based competitor
  would use today, and \citet{HochbaumQueyranne2003} place the closure problem
  and isotonic regression in one family.
\item[The ratio objective and the decomposition.] The normalization is
  \citet{CharnesCooper1962}; the parametric reading is
  \citet{Dinkelbach1967}; the decomposition into blocks of decreasing ratio is
  \citet{Sidney1975} and \citet{Lawler1978}; parametric flow over the whole
  ratio range is \citet{GalloEtAl1989} and \citet{Eppstein2018}, and the
  divide-and-conquer scheme is \citet{EisnerSeverance1976}. The union
  convention that makes a block canonical is the maximal source side of a
  minimum cut \citep{PicardQueyranne1980}.
\item[Forests carrying branch aggregates.] This is the oldest and closest
  antecedent. The open-pit design algorithm of \citet{LerchsGrossmann1965}
  maintains a \emph{normalized tree}, carries the mass of each branch,
  classifies an arc by the sign of the mass it supports, and pivots by adding
  one arc and removing a branch's root arc; \citet{Hochbaum2001} sets this out
  and recovers a flow from the tree in linear time, and also gives a parametric
  variant with the cost of a single run. Substituting $b_i=p_i-\rho w_i$ turns
  the strong/weak test into the sign of the reduced costs derived in
  \cref{sec:directions}.
\item[Recent simplex work on the same shapes.] \citet{Yang2026} develops the
  revised-simplex algebra on rooted spanning forests for
  budget-constrained total-variation problems, which by reading (c) of
  \cref{sec:class} is the same structural family as \eqref{eq:normalized}.
  \citet{Wuille2025} states \eqref{eq:normalized} itself, with $n-1$ nonbasic
  variables, a spanning forest of
  tight dependency slacks, at most one free vertex variable per component and a
  single positive component; \citet{MollakhaniEtAl2026} formalize that as a
  spanning-forest linearization algorithm. Simplex methods on an incidence
  matrix with one side row are classical
  \citep{AhujaEtAl1993,Schrijver2003,HolzhauserEtAl2017}.
\end{description}

What this tutorial develops is therefore a \emph{specialization}: the ordinary
primal simplex carried through exactly on \eqref{eq:normalized}, with the closed
forms that specialization produces. It is not a new class of basis and not a
new paradigm.

\subsection{Cycles and limits of the model}
If the input digraph contains a directed cycle, the corresponding inequalities
force equality of the variables on that cycle. Contract each strongly
connected component, sum its profits and weights, and retain a map to expand
the final solution. The resulting graph is a DAG with positive aggregated
weights. Thus the DAG assumption is a preprocessing convention for general
directed precedence graphs. In contrast, zero or negative weights, several
independent capacity rows, heterogeneous upper bounds, and inequalities of
the form $x_i\le\gamma_{ij}x_j$ require additional analysis. They are not
included in the statements of this draft.
